% demo_petritikz.tex
% ─────────────────────────────────────────────────────────────────
%  Compilation :  pdflatex --shell-escape demo_petritikz.tex
%  (Windows et Linux/Mac, pas de flag supplementaire)
%
%  Fichiers requis dans le meme dossier :
%    petritikz.sty    styles TikZ + commande \petrinet
%    petritikz.py     generateur Python
%
%  Utilisation :
%    \petrinet{--pre 1,0,0,0;0,1,1,0 --post 0,1,1,0;1,0,0,1
%              --places libre,produit,vide,plein
%              --trans  produire,consommer
%              --marking 1,0,1,0  --layout circle}
%
%  Format des matrices :  lignes separees par ;  valeurs par ,
% ─────────────────────────────────────────────────────────────────

\documentclass[a4paper,11pt]{article}
\usepackage[T1]{fontenc}
\usepackage[utf8]{inputenc}
\usepackage{petritikz}
\usepackage{amsmath,amssymb}
\usepackage{booktabs,array}
\usepackage{xcolor}
\usepackage{geometry}
\usepackage{fancyvrb}
\usepackage{tcolorbox}
\tcbuselibrary{skins,breakable}
\geometry{margin=2.5cm}

\newtcolorbox{codebox}{
  colback=green!4, colframe=green!40!black,
  fonttitle=\bfseries\ttfamily\small,
  title={Code dans le .tex},
  left=5pt, right=5pt, top=3pt, bottom=3pt
}
\newtcolorbox{resbox}{
  colback=blue!4, colframe=blue!30,
  fonttitle=\bfseries\small, title={Resultat},
  left=5pt, right=5pt, top=3pt, bottom=3pt
}

\newcolumntype{C}[1]{>{\centering\arraybackslash}p{#1}}

\title{\textbf{PetriTikz}\\[4pt]
  {\large La commande \texttt{\textbackslash petrinet\{\}} appelle Python}}
\author{PetriTikz}
\date{\today}

\begin{document}
\maketitle

% ═══════════════════════════════════════════════════════════════════
\section{Reseau producteur--consommateur}
% ═══════════════════════════════════════════════════════════════════

\subsection{Schema généré automatiquement}

\begin{center}
\petrinet{--pre 0,1,0;1,0,0;0,0,1;0,1,0 --post 1,0,0;0,1,0;0,1,0;0,0,1 --marking 1,0,1,0 --layout hierarchical --struct V}
\end{center}

\begin{center}
\petrinet{--pre 1,0,0,0,0,0;0,1,0,0,0,0;0,0,1,0,0,0;0,1,0,0,1,0;0,0,0,1,0,0;0,0,0,0,1,0;0,0,0,0,0,1 --post 0,1,0,0,0,0;1,0,0,0,0,0;0,1,0,0,0,0;0,0,1,0,0,1;0,0,0,0,1,0;0,0,0,1,0,0;0,0,0,0,1,0 --marking 1,0,0,1,1,0,0 --layout hierarchical --struct V}
\end{center}


\end{document}